\documentclass[../../../include/open-logic-section]{subfiles}

\begin{document}

\olfileid[id]{sth}{ordinals}{wo}
\olsection{Pengurutan Baik}

Gagasan dasarnya adalah sebagai berikut:

\begin{defn}
Relasi~$<$ \emph{mengurutkan $A$ dengan baik} jika dan hanya jika relasi itu
memenuhi kedua syarat berikut:
\begin{enumerate}
	\item $<$ terhubung, yakni untuk semua $a,b\in A$, berlaku $a<b$ atau
	$a=b$ atau $b<a$;
	\item setiap himpunan bagian tak kosong dari~$A$ mempunyai suatu
	!!{element} yang minimal menurut~$<$, yakni jika
	$\emptyset\neq X\subseteq A$, maka
	$(\exists m\in X)(\forall z\in X)z\nless m$.
\end{enumerate}
\end{defn}

Mudah dilihat bahwa ketiga contoh yang baru saja kita pertimbangkan memang
merupakan relasi pengurutan baik.

\begin{prob}
\Olref[sth][ordinals][idea]{sec} menyajikan tiga contoh pengurutan pada
bilangan asli. Periksa bahwa masing-masing merupakan pengurutan baik.
\end{prob}

Berikut beberapa pengamatan elementer tetapi sangat penting mengenai
pengurutan baik.

\begin{prop}\ollabel{wo:strictorder}
Jika $<$ mengurutkan $A$ dengan baik, maka setiap himpunan bagian tak kosong
dari~$A$ mempunyai anggota terkecil yang unik menurut~$<$, dan $<$ bersifat
irefleksif, asimetris, serta transitif.
\end{prop}

\begin{proof}
Jika $X$ merupakan himpunan bagian tak kosong dari~$A$, maka $X$ mempunyai
suatu !!{element}~$m$ yang minimal menurut~$<$, yakni
$(\forall z\in X)z\nless m$. Karena $<$ terhubung, berlaku
$(\forall z\in X)m\leq z$. Jadi, $m$ merupakan !!{element} terkecil dari~$X$
menurut~$<$.

Untuk irefleksivitas, tetapkan $a\in A$; !!^{element} terkecil dari $\{a\}$
menurut~$<$ ialah $a$, sehingga $a\nless a$. Untuk transitivitas, andaikan
$a<b<c$. Jika $a\nless c$, keterhubungan memberikan $a=c$ atau $c<a$.
  Dalam kasus pertama, $\{a,b\}$, dan dalam kasus kedua, $\{a,b,c\}$, tidak
  mempunyai suatu !!{element} minimal. Keduanya mustahil. Jadi,
$a<c$. Asimetri mengikuti irefleksivitas dan transitivitas.
\end{proof}

\begin{prop}\ollabel{propwoinduction}
Jika $<$ mengurutkan $A$ dengan baik, maka untuk sebarang formula~$\phi(x)$:
\[
	\text{jika }(\forall a \in A)((\forall b < a)\phi(b) \lif 
		\phi(a))\text{, maka }(\forall a \in A)\phi(a).
\]
\end{prop}

\begin{proof}
Kita membuktikan kontraposisinya. Andaikan
$\lnot(\forall a\in A)\phi(a)$, yakni
$X=\Setabs{x\in A}{\lnot\phi(x)}\neq\emptyset$. Maka $X$ mempunyai suatu
!!{element} minimal~$a$ menurut~$<$. Jadi, $(\forall b<a)\phi(b)$, tetapi
$\lnot\phi(a)$.
\end{proof}\noindent
Sifat terakhir ini seharusnya mengingatkan Anda pada prinsip induksi kuat
pada bilangan asli, yakni: jika
$(\forall n\in\omega)((\forall m<n)\phi(m)\lif\phi(n))$, maka
$(\forall n\in\omega)\phi(n)$. Sifat inilah yang menjadikan pengurutan baik
sebagai gagasan yang sangat \emph{tangguh}.\footnote{Pengingat: semua formula
dapat mempunyai parameter (kecuali dinyatakan sebaliknya secara eksplisit).}

\end{document}
